\section{Conclusion}
Dans ce travail, nous avons étudié un système de lévitation magnétique, reconnu pour son comportement non linéaire et instable en boucle ouverte. Nous avons d’abord établi le modèle d’état non linéaire du système, puis procédé à sa linéarisation autour d’un point d’équilibre afin de faciliter l’analyse.

L’étude du modèle linéarisé a permis de déterminer la fonction de transfert ainsi que les pôles du système. L’identification des dynamiques dominantes a mené à la simplification du modèle en un système d’ordre 2, ce qui a facilité la conception et l’analyse du correcteur.

Par la suite, une loi de commande de type PD a été analysée et ajustée à partir de données expérimentales. Les résultats obtenus ont montré une bonne concordance entre le modèle identifié et la réponse réelle du système, validant ainsi l’approche utilisée.

Finalement, les simulations du modèle non linéaire en boucle fermée ont permis d’observer le comportement du système face à différentes consignes. Bien que la sortie suive correctement la tendance de la consigne, un écart persistant a été observé, suggérant l’ajout d’une action intégrale pour améliorer la précision du suivi.

En conclusion, ce travail met en évidence l’importance de la modélisation, de la linéarisation et de la conception de correcteurs dans l’analyse et le contrôle de systèmes non linéaires complexes.\label{sec:conclusion}

